\documentclass[11pt]{article}

% ---------------------------------------------------------------------------
% Environmental effects on meiotic recombination landscapes
% Plain-language working draft
% Compiles with pdfLaTeX. Keep model_illustration_revised.pdf in this folder.
% ---------------------------------------------------------------------------

\usepackage[margin=1in]{geometry}
\usepackage{amsmath,amssymb,amsthm,mathtools}
\usepackage{graphicx}
\usepackage{booktabs,tabularx,array}
\usepackage{microtype}
\usepackage[round,authoryear]{natbib}
\usepackage[hidelinks,hypertexnames=false]{hyperref}
\usepackage{xcolor}
\usepackage{enumitem}
\usepackage{float}

\renewcommand{\topfraction}{0.9}
\renewcommand{\textfraction}{0.08}
\renewcommand{\floatpagefraction}{0.8}
\setlist{nosep,leftmargin=*}

\newtheorem{assumption}{Assumption}
\newtheorem{proposition}{Proposition}
\newtheorem{prediction}{Prediction}
\theoremstyle{remark}
\newtheorem*{remark}{Explanation}
\newtheorem*{scope}{Scope}

\newcommand{\dd}{\,\mathrm{d}}
\newcommand{\E}{\mathbb{E}}
\newcommand{\avg}[1]{\langle #1 \rangle}
\newcommand{\Rbar}{\overline{R}}
\newcommand{\pbar}{\overline{p}}
\newcommand{\FST}{F_{\mathrm{ST}}}
\newcommand{\dxy}{d_{xy}}
\newcommand{\PRDM}{\textsc{prdm9}}
\newcommand{\SPO}{\textsc{spo11}}

\title{\bfseries Environmental effects on meiotic recombination maps\\[6pt]
{\large\normalfont A model linking meiotic chromatin to diversity and relative differentiation}}

\author{
Emily C.\ Giles\,$^{1}$, \;
Pablo Saenz-Agudelo\,$^{1}$, \;
Daniel Ortiz-Barrientos\,$^{2,\ast}$
}
\date{\today \\[6pt]
\small
$^{1}$Cawthron Institute, 98 Halifax Street East, Nelson 7010,
New Zealand (postal: Private Bag 2, Nelson 7040) \\[1pt]
$^{2}$School of the Environment and ARC Centre of Excellence for Plant Success in
Nature and Agriculture, \\
The University of Queensland, St Lucia QLD 4072, Australia \\[3pt]
\small $^{\ast}$Corresponding author: \texttt{d.ortizbarrientos@uq.edu.au}}

\begin{document}
\maketitle

\begin{abstract}
\noindent
Environmental conditions can change both the number and the genomic positions of
meiotic crossovers. Population-genetic models usually treat the recombination map
as fixed, and experimental studies often combine these two changes. We separate
them. In environment $E$, we write the map as $r_E(x)=R_Ep_E(x)$, where $R_E$ is
total map length and $p_E(x)$ is the distribution of crossover positions. A local
placement weight $q_E(x)$ determines $p_E(x)$. Meiotic chromatin accessibility
contributes to this weight together with sequence context, chromosome structure,
axis organisation, and crossover maturation. For two environments, the map change
is $\Delta r=\Delta R\,\pbar+\Rbar\,\Delta p$. The first term measures the change
in total map length. The second measures the redistribution of crossovers along the
chromosome and integrates to zero. A region gains relative crossover probability
when its increase in placement weight exceeds its expected share of the genome-wide
increase. The result also applies to regions with zero placement weight at
baseline. We then model background selection, in which selection against deleterious alleles
also removes linked neutral variants. In a local approximation, diversity rises
where recombination rises. Relative
differentiation falls in those regions when absolute divergence changes little. In
the full model, diversity depends on recombination along the chromosome between
neutral and selected sites, so the sign depends on changes across the surrounding
map. These population-genetic effects accumulate over many generations. They can arise when the environment acts directly on meiosis in each generation
or when an inherited state persists after exposure. The model specifies experiments and comparative tests that link meiotic chromatin,
recombination-map change, and population-genomic variation.
\end{abstract}

\section{Introduction}

Crossovers are unevenly distributed along chromosomes. In budding yeast and
plants, open chromatin, nucleosome organisation, histone modifications, DNA
methylation, chromosome-axis proteins, sequence polymorphism, and structural
context all contribute to their positions
\citep{pan2011,choi2013,choi2018,underwood2018,fernandes2024,son2025}.
Experiments that alter chromatin states also alter meiotic recombination, which
supports a causal role for chromatin
\citep{underwood2018,fernandes2024,son2025,szymanska2026}.

The recombination map also responds to the environment. Temperature can change
total crossover number, crossover positions, or both. The response varies among
chromosomes, sexes, genotypes, and temperature treatments
\citep{lloyd2018,modliszewski2018,destorme2020,coulton2020}. An analysis must distinguish total map length from the spatial distribution of
crossovers.

Recombination also affects neutral genetic variation. Selection at one site can
remove neutral variants at linked sites. Recombination weakens this association.
Background selection and selective sweeps reduce more neutral diversity
in regions with low recombination
\citep{begun1992,charlesworth1993,hudson1995,nordborg1996,coop2012}. The same
process can raise relative differentiation by reducing diversity within
populations, even when absolute divergence remains near its background value
\citep{charlesworth1998,cruickshank2014}. Most models treat the recombination map
as fixed. Theory on evolving maps shows that map history can change the strength
and spatial distribution of background selection \citep{booker2022}. Recombination
maps can also differ between recently diverged ecotypes \citep{talbi2025}.

We develop one model for these observations. We first separate total map length
from the spatial distribution of crossovers. We then derive how an environmental
change in local crossover placement alters that distribution. Next, we calculate
the effects on diversity under local and nonlocal models of linked selection. We
also derive the corresponding change in relative differentiation. Each equation
states the conditions required for its use.

The map response and its population-genetic consequences can occur without
epigenetic inheritance. The environment may act directly on cells entering meiosis,
or the same environment may induce a similar state in each generation. Inheritance becomes a separate hypothesis when
an altered state persists after the inducing environment has been removed.

\section{From meiotic state to the recombination map}
\label{sec:model}

\subsection{A model for crossover placement}

Let $x\in[0,L]$ denote physical position along a chromosome and let $E$ denote an
environment or treatment. Define a nonnegative, integrable crossover-placement
weight
\begin{equation}
q_E(x)=f\!\left(a_E(x),\mathbf{z}_E(x)\right),
\qquad q_E(x)\ge 0,
\label{eq:q}
\end{equation}
where $a_E(x)$ is meiotic chromatin accessibility. The vector $\mathbf{z}_E(x)$
contains other determinants of crossover placement, including chromatin marks,
DNA methylation, sequence and structural variation, chromosome-axis organisation,
synapsis, interference, sex, and the probability that a double-strand break
matures into a crossover. The weight $q_E(x)$ measures local crossover propensity
before normalisation. We leave the function $f$ unspecified. The empirical
question is whether accessibility explains part of $q_E$ after the other
determinants have been controlled.

Let
\begin{equation}
Q_E=\int_0^L q_E(x)\dd x,
\qquad
p_E(x)=\frac{q_E(x)}{Q_E},
\qquad
\int_0^L p_E(x)\dd x=1.
\label{eq:p}
\end{equation}
The density $p_E(x)$ gives the spatial distribution of crossover positions,
conditional on a crossover occurring. Let $R_E$ be total map length in environment
$E$. The realised map is
\begin{equation}
r_E(x)=R_Ep_E(x),
\qquad
\int_0^L r_E(x)\dd x=R_E.
\label{eq:r}
\end{equation}
This factorisation allows total map length and spatial distribution to vary
independently.

\begin{scope}
The model applies where meiotic chromatin explains a measurable part of crossover
placement. A functional \PRDM{} system may be present or absent. In
\PRDM-directed systems, $q_E$ can include motif binding and the chromatin changes
recruited by \PRDM. In systems without functional \PRDM-directed placement,
promoters, CpG islands, and other chromatin-associated features often explain more
of the placement pattern \citep{baker2017,talbi2025}. The model should be tested by
estimating the contribution of meiotic chromatin after accounting for the other
terms in $\mathbf{z}_E$.
\end{scope}

\subsection{Total map length and spatial distribution}

Compare a baseline environment $E=0$ with a second environment $E=1$. Write
\[
\Delta r=r_1-r_0,
\quad
\Delta R=R_1-R_0,
\quad
\Delta p=p_1-p_0,
\]
and define the midpoints
\[
\Rbar=\frac{R_1+R_0}{2},
\qquad
\pbar=\frac{p_1+p_0}{2}.
\]

\begin{proposition}[exact total--shape decomposition]
\label{prop:decomp}
For any two recombination maps of the form \eqref{eq:r},
\begin{equation}
\boxed{\;
\Delta r(x)=\underbrace{\Delta R\,\pbar(x)}_{\text{total-map component}}
+\underbrace{\Rbar\,\Delta p(x)}_{\text{shape component}}.
\;}
\label{eq:decomp}
\end{equation}
The two components satisfy
\begin{equation}
\int_0^L \Delta R\,\pbar(x)\dd x=\Delta R,
\qquad
\int_0^L \Rbar\,\Delta p(x)\dd x=0.
\label{eq:decomp_integrals}
\end{equation}
\end{proposition}

\begin{proof}
Expanding the right-hand side gives
\[
\frac{(R_1-R_0)(p_1+p_0)+(R_1+R_0)(p_1-p_0)}{2}
=R_1p_1-R_0p_0.
\]
The integral statements follow from $\int p_0=\int p_1=1$.
\end{proof}

\begin{remark}
The zero integral of the shape component follows from normalisation. It describes
redistribution along the chromosome. Experiments must estimate $\Delta R$
separately to determine whether total recombination has changed.
\end{remark}

We measure the amount of spatial redistribution in map units as
\begin{equation}
M_{\mathrm{shape}}
=\frac{\Rbar}{2}\int_0^L |\Delta p(x)|\dd x.
\label{eq:mshape}
\end{equation}
The factor $1/2$ avoids counting the same redistributed map length twice. Let $M_{\mathrm{total}}=|\Delta R|$. When either component is
nonzero, define
\begin{equation}
\chi
=\frac{M_{\mathrm{shape}}}
{M_{\mathrm{shape}}+M_{\mathrm{total}}}.
\label{eq:chi}
\end{equation}
Values near one indicate a response dominated by spatial redistribution. Values
near zero indicate a response dominated by total map-length change. We use $\chi$
as a descriptive summary of the observed response.

\subsection{Change in local crossover placement}
\label{sec:shape}

Let $q_1=q_0+\Delta q$ and $Q_1=Q_0+\Delta Q$, where
$\Delta Q=\int_0^L\Delta q(x)\dd x$. Write $p=p_0$ for the baseline distribution.

\begin{proposition}[exact redistribution of normalised placement]
\label{prop:placement}
For any admissible perturbation with $Q_1>0$,
\begin{equation}
\boxed{\;
\Delta p(x)
=\frac{\Delta q(x)-p(x)\Delta Q}{Q_1}.
\;}
\label{eq:dp_general}
\end{equation}
Consequently,
\begin{equation}
\Delta p(x)>0
\quad\Longleftrightarrow\quad
\Delta q(x)>p(x)\Delta Q.
\label{eq:dp_sign_general}
\end{equation}
\end{proposition}

\begin{proof}
Direct subtraction gives
\[
\frac{q+\Delta q}{Q+\Delta Q}-\frac{q}{Q}
=\frac{Q\Delta q-q\Delta Q}{Q(Q+\Delta Q)}
=\frac{\Delta q-p\Delta Q}{Q+\Delta Q}.
\]
The denominator is positive, so the numerator determines the sign.
\end{proof}

Equation~\eqref{eq:dp_general} remains defined when $q(x)=0$. If an inactive
region acquires positive placement weight, then
$\Delta p(x)=\Delta q(x)/Q_1>0$. The equation also covers the appearance of new
crossover-active regions.

Where $q(x)>0$, define the fractional change $h(x)=\Delta q(x)/q(x)$ and its
baseline placement-weighted mean
\begin{equation}
\avg{h}_p=\int_0^L p(x)h(x)\dd x=\frac{\Delta Q}{Q}.
\end{equation}
Then Eq.~\eqref{eq:dp_general} becomes
\begin{equation}
\boxed{\;
\Delta p(x)
=p(x)\frac{h(x)-\avg{h}_p}{1+\avg{h}_p}.
\;}
\label{eq:dp_fractional}
\end{equation}
A region gains relative crossover probability when its placement weight increases
by more than the placement-weighted genomic average. A uniform multiplicative
change in $q$ changes $Q$ and leaves $p$ unchanged.

Chromatin accessibility enters this result as one component of $q_E$. If
accessibility changes while the other arguments of $f$ remain fixed, then a small
perturbation at locations with $a>0$ gives
\begin{equation}
h(x)
\approx
\varepsilon_a(x)g_a(x),
\qquad
\varepsilon_a(x)=\frac{\partial\log f}{\partial\log a},
\qquad
g_a(x)=\frac{\Delta a(x)}{a(x)}.
\label{eq:elasticity}
\end{equation}
If $f(a)=a$, then $h=g_a$ exactly. More general functions allow accessibility to
change double-strand-break initiation while other processes control crossover
maturation. Models for $q_E$ should include all available predictors of crossover placement.

\begin{figure}[t]
\centering
\includegraphics[width=0.98\textwidth]{model_illustration_plain.pdf}
\caption{An example of an environmental change in the recombination map.
\textbf{(a)} The environment changes the local crossover-placement weight
$q_E(x)$, which includes meiotic chromatin and other determinants.
\textbf{(b)} Total map length and the spatial distribution can change together.
\textbf{(c)} Equation~\eqref{eq:decomp} separates the observed difference into a
total component, which integrates to $\Delta R$, and a shape component, which
integrates to zero. \textbf{(d)} The local background-selection approximation
predicts higher diversity where the local recombination rate rises. Relative
differentiation falls in the same region when $\dxy$ changes little. The nonlocal
response in Eq.~\eqref{eq:dlogpi_nonlocal} averages map changes over the paths
between neutral and selected sites.}
\label{fig:model}
\end{figure}

\section{Consequences for linked selection}
\label{sec:diversity}

\subsection{Background selection across the recombination map}

Background selection reduces neutral diversity when selection removes deleterious
alleles and the neutral variants linked to them. Recombination separates neutral
and selected sites, so the effect declines with genetic distance. Let $u(z)$ be
the deleterious mutation-rate density at position $z$, and let $s(z)>0$ be the
selection coefficient. Under map $r$, the genetic distance between positions $x$
and $z$ is
\begin{equation}
d_r(x,z)
=\int_{\min(x,z)}^{\max(x,z)} r(y)\dd y.
\label{eq:distance}
\end{equation}
We use the following standard approximation for the background-selection exponent:
\begin{equation}
E(x;r)
=\int_0^L u(z)
\frac{s(z)}{\left[s(z)+d_r(x,z)\right]^2}
\dd z,
\qquad
\pi(x;r)=\pi_{\mathrm{neu}}(x)e^{-E(x;r)}.
\label{eq:bgs}
\end{equation}
Here $\pi_{\mathrm{neu}}(x)$ is neutral diversity in the absence of linked
selection. The contribution from a selected site decreases as its genetic
distance from $x$ increases.

For a small map perturbation $\delta r$, the change in genetic distance is
\begin{equation}
\delta d_r(x,z)
=\int_{\min(x,z)}^{\max(x,z)}\delta r(y)\dd y.
\label{eq:ddistance}
\end{equation}
Differentiating Eq.~\eqref{eq:bgs} gives
\begin{equation}
\boxed{\;
\delta\log\pi(x)
=2\int_0^L u(z)
\frac{s(z)}{\left[s(z)+d_r(x,z)\right]^3}
\left[
\int_{\min(x,z)}^{\max(x,z)}\delta r(y)\dd y
\right]\dd z.
\;}
\label{eq:dlogpi_nonlocal}
\end{equation}
The ordered limits in Eq.~\eqref{eq:ddistance} give the same positive genetic
distance for selected sites on either side of $x$.

Equation~\eqref{eq:dlogpi_nonlocal} can also be written as
\begin{equation}
\delta\log\pi(x)=\int_0^L W_x(y)\,\delta r(y)\dd y,
\qquad W_x(y)\ge 0,
\label{eq:kernel}
\end{equation}
where $W_x(y)$ weights positions on the chromosomal segments between $x$ and
selected sites. Thus the change in diversity at $x$ is a weighted average of map
changes across those segments. Changes away from $x$ can reinforce or offset the
local change. A pointwise sign rule requires the local approximation below.

\subsection{Local approximation}

Suppose $u$, $s$, and $r$ vary slowly over the linkage scale that contributes most
to Eq.~\eqref{eq:bgs}. Also suppose that the site of interest lies far enough from a chromosome
boundary. Then
\begin{equation}
E(x;r)\approx\frac{\kappa u(x)}{r(x)},
\qquad \kappa=2
\label{eq:bgs_local}
\end{equation}
for the kernel in Eq.~\eqref{eq:bgs} on a locally uniform, effectively infinite
chromosome. Let $E_0(x)=\kappa u(x)/r_0(x)$. For a finite local map change,
\begin{equation}
\boxed{\;
\log\frac{\pi_1(x)}{\pi_0(x)}
=E_0(x)\left[1-\frac{r_0(x)}{r_1(x)}\right].
\;}
\label{eq:pi_ratio_general}
\end{equation}
Within this approximation, $r_1(x)>r_0(x)$ gives $\pi_1(x)>\pi_0(x)$.

For $q_0(x)>0$, let $h=\Delta q/q_0$ and let
$\eta_R=\Delta R/R_0$. From Eqs.~\eqref{eq:r} and
\eqref{eq:dp_fractional},
\begin{equation}
\frac{r_1(x)}{r_0(x)}
=(1+\eta_R)
\frac{1+h(x)}{1+\avg{h}_p}.
\label{eq:r_ratio}
\end{equation}
Substitution into Eq.~\eqref{eq:pi_ratio_general} gives
\begin{equation}
\boxed{\;
\log\frac{\pi_1(x)}{\pi_0(x)}
=E_0(x)
\left[
1-\frac{1+\avg{h}_p}
{(1+\eta_R)(1+h(x))}
\right].
\;}
\label{eq:pi_ratio_h}
\end{equation}
For small changes,
\begin{equation}
\frac{\Delta\pi(x)}{\pi(x)}
\approx
E_0(x)\left[\eta_R+h(x)-\avg{h}_p\right].
\label{eq:pi_firstorder}
\end{equation}
When $\eta_R=0$, the local diversity sign follows
$h-\avg{h}_p$. When total map length changes, the total and spatial components may
reinforce or oppose each other.

Both the local and nonlocal formulations give two general results. Regions with
strong existing linked selection show larger fractional responses to a map change.
Genome-wide diversity can also change when $\Delta p$ integrates to zero, because
the response is nonlinear and selected sites are unevenly distributed.

\subsection{Timescale and map history}
\label{sec:timescale}

A recombination map can respond within one meiosis. Nucleotide diversity changes
as genealogies turn over across many generations. If a map remains stable long
enough, Eqs.~\eqref{eq:bgs}--\eqref{eq:pi_ratio_h} describe the resulting long-term
values. If the map switches before genealogies relax, current diversity depends
on the sequence of past maps \citep{booker2022}. Forward or coalescent simulations
can use that sequence directly. A one-generation experiment can estimate
$\Delta a$, $\Delta q$, $\Delta R$, and $\Delta p$. Tests of $\Delta\pi$ require
recurrent exposure, a persistent state, or comparisons among populations with
different long-term map histories.

\section{Consequences for relative differentiation}
\label{sec:divergence}

For a Hudson-type sequence measure, write relative differentiation as
\begin{equation}
\FST(x)=1-\frac{\pi_w(x)}{\dxy(x)},
\label{eq:fst_def}
\end{equation}
where $\pi_w$ is diversity within populations and $\dxy$ is mean pairwise
divergence between populations. For two map regimes,
\begin{equation}
\boxed{\;
1-F_{\mathrm{ST},1}(x)
=\left[1-F_{\mathrm{ST},0}(x)\right]
\frac{\pi_{w,1}(x)}{\pi_{w,0}(x)}
\frac{d_{xy,0}(x)}{d_{xy,1}(x)}.
\;}
\label{eq:fst_finite}
\end{equation}
Its first-order form is
\begin{equation}
\Delta\FST(x)
\approx
\left[1-\FST(x)\right]
\left[
\frac{\Delta\dxy(x)}{\dxy(x)}
-
\frac{\Delta\pi_w(x)}{\pi_w(x)}
\right].
\label{eq:fst_firstorder}
\end{equation}

Consider a map change that occurs after population separation, leaves the neutral
mutation rate nearly unchanged, and acts mainly through within-population linked
selection. In this regime, $d_{xy,1}\approx d_{xy,0}$. The local
background-selection approximation then gives
\begin{equation}
\boxed{\;
F_{\mathrm{ST},1}(x)
=1-\left[1-F_{\mathrm{ST},0}(x)\right]
\exp\!\left\{
E_0(x)\left[1-\frac{r_0(x)}{r_1(x)}\right]
\right\}.
\;}
\label{eq:fst_local}
\end{equation}
Under these conditions,
\begin{equation}
\operatorname{sign}\!\left(F_{\mathrm{ST},1}-F_{\mathrm{ST},0}\right)
=-\operatorname{sign}(r_1-r_0).
\label{eq:fst_sign}
\end{equation}
When $\eta_R=0$, the sign is opposite to $h-\avg{h}_p$. This sign rule applies to
the stated local linked-selection regime with approximately stable $\dxy$.

The analysis should estimate $\dxy$. Ancestral polymorphism, mutation-rate
variation, gene flow, introgression, ancestral linked selection, and structural
variation can all change it. High $\FST$ together with low $\pi_w$ and $\dxy$ at or below its genomic
background is consistent with linked selection in a low-recombination region.
Elevated $\dxy$ can arise from restricted gene flow or divergent selection, among
other causes. The analysis should combine direct recombination maps, $\pi_w$, linkage
disequilibrium, $\FST$, $\dxy$, mutation-rate proxies, and demographic models.

Replicate populations exposed to similar environments may develop similar $\FST$
profiles even when selection acts on different loci. Evidence for the proposed
mechanism requires a similar meiotic chromatin response, a similar $\Delta p$, and
an $\FST$ response mediated mainly through $\pi_w$. Other possible causes include
parallel selection, structural variants, and shared ancestral recombination maps.

\section{Direct exposure, repeated induction, and inheritance}
\label{sec:inheritance}

Three processes can produce recurrent differences in recombination maps.

\paragraph{Direct parental exposure.}
The environment reaches cells that undergo meiosis and changes $q_E$, $R_E$, or
both. The altered gametes record a parental environmental effect.

\paragraph{Repeated environmental induction.}
The meiotic state resets each generation, and the same environment induces a
similar map in each generation. Population-genomic effects can accumulate without
persistence of the chromatin state between generations.

\paragraph{Inherited state.}
The inducing environment ends, and the altered meiotic state persists in one or
more unexposed generations. A transgenerational inheritance claim requires
persistence in the relevant unexposed generations.

A minimal model separates direct and inherited components. Let $E_t$ be an
environmental cue experienced by generation $t$, let $a_0(x)$ be the baseline
meiotic accessibility profile, and write
\begin{align}
a_t(x)&=a_0(x)+\beta(x)E_t+m_t(x),
\label{eq:direct_memory_a}\\
m_{t+1}(x)&=\rho(x)m_t(x)+\gamma(x)E_t+\xi_t(x),
\qquad 0\le \rho(x)\le 1.
\label{eq:direct_memory_m}
\end{align}
The coefficient $\beta$ gives the direct response within a generation. The state
$m_t$ is the transmitted component. The coefficient $\gamma$ measures how exposure
changes that component, and $\rho$ measures its persistence after $E_t$ returns to
zero. Repeated induction can operate with $m_t=0$. Evidence for inheritance
requires $m_t\ne0$ in the relevant unexposed descendants. The number of unexposed
generations depends on the organism's germline and life cycle.

We use the following terms. Repeated environmental induction is an
environmental input to the recombination map. Persistent $m_t$ is an inherited
epigenetic input. We reserve the term feedback for systems in which the resulting
genetic variation later changes the chromatin response or the interaction with the
environment.

\section{Tests of the model}
\label{sec:predictions}

The model contains a sequence of testable links. Evidence for population-genomic
effects depends on support for the earlier meiotic and recombination links.



\subsection{Estimate crossover number and position separately}

The primary experiment should retain crossover data from individual meioses. Fit
total crossover number or total map length by bivalent, meiocyte, sex, genotype,
and treatment. This estimates $R_E$ and the extent of crossover homeostasis. Then
condition on the observed number of crossovers and model their positions along the
chromosome. This estimates $p_E$ and tests the placement model. A point-process
model can include crossover interference. A multinomial or Dirichlet-multinomial
model may be adequate for coarse genomic bins. Binwise correlations can mislead
because $p_E$ is compositional: an increase in one region produces relative
decreases elsewhere. Conditional likelihoods or log-ratio methods account for
this constraint.

Collect the chromatin data at the meiotic stage that produces the mapped gametes.
Use somatic accessibility after establishing its relation to the meiotic state. Analyse male and female maps separately when possible. Viability
selection in gametes or zygotes can distort offspring-based crossover maps. Direct
gamete assays and reciprocal designs can test for this bias.

\begin{table}[H]
\centering
\small
\begin{tabularx}{\textwidth}{>{\raggedright\arraybackslash}p{0.08\textwidth}
>{\raggedright\arraybackslash}p{0.22\textwidth}
>{\raggedright\arraybackslash}X
>{\raggedright\arraybackslash}p{0.24\textwidth}}
\toprule
Link & Quantity & Required design and expected result & Evidence against the proposed link \\
\midrule
P0 & Baseline coupling & In meiocytes, estimate whether accessibility and other chromatin variables explain crossover placement after controlling for sequence, structure, gene density, chromosome position, sex, and genotype. & Accessibility adds no predictive information to crossover placement in the study system. \\
\addlinespace
P1 & Environmental mechanism & Measure meiotic chromatin and crossover locations in matched control and treatment samples. Estimate $\Delta q$ from the molecular predictors and test whether it predicts $\Delta p$ through Eq.~\eqref{eq:dp_general}. & The observed $\Delta p$ is repeatedly unrelated to the measured placement-state change after relevant covariates and meiotic timing are controlled. \\
\addlinespace
P2 & Total and spatial change & Estimate $R_E$ and $p_E$ separately, then report $\Delta R$, $M_{\mathrm{shape}}$, and $\chi$. & A large $\Delta R$ with little $\Delta p$ supports a change in total map length and gives little evidence for spatial redistribution. \\
\addlinespace
P3 & Long-run diversity & In replicated populations with recurrent exposure or a persistent map shift, test whether measured maps predict $\pi$ and LD under forward or coalescent models that include demography and linked selection. & A fixed-map model fits as well as the measured map history, or the observed response differs from the prediction based on that history. \\
\addlinespace
P4 & Relative differentiation & Jointly model $\pi_w$, LD, $\FST$, and $\dxy$. In the stated linked-selection regime, changes in $\FST$ should arise mainly through changes in $\pi_w$. & Mutation-rate variation, structural variants, gene-flow barriers, or selection at the loci under study explain the data better, and direct map data give no support for the proposed mechanism. \\
\addlinespace
P5 & Inheritance & Remove the cue and assay meiotic chromatin and recombination in the required unexposed descendants. Estimate persistence in Eqs.~\eqref{eq:direct_memory_a}--\eqref{eq:direct_memory_m}. & The state returns to baseline in all unexposed descendants. The result is consistent with direct or repeated induction and excludes persistent memory. \\
\bottomrule
\end{tabularx}
\caption{Tests of the proposed mechanism. The zero integral of $\Delta p$ follows
from normalisation. Empirical tests concern changes in map shape, the contribution
of meiotic chromatin to those changes, and the duration of the altered map.}
\label{tab:predictions}
\end{table}

\subsection{Predict population-genomic effects}

After estimating $\Delta R$ and $\Delta p$, use the measured map in forward or
coalescent simulations. The model can include deleterious mutations, selective
sweeps, demography, structural variants, and the frequency of environmental
exposure. Simulations are needed because Eqs.~\eqref{eq:pi_ratio_h} and
\eqref{eq:fst_local} describe long-term local values. They should predict the spatial and
temporal responses of $\pi$, LD, $\FST$, and $\dxy$ before comparisons with natural
populations.

Comparative tests require ecological replication. Suitable designs include
independent population pairs, replicate colonisations, and populations distributed
along an environmental gradient. The analysis should include mutation rate, gene
density, mappability, GC-biased gene conversion, centromeres, structural variants,
sex-specific maps, and demographic history. LD-based recombination maps depend on
demography and selection. Pedigree, gamete, and experimental maps provide more
direct tests of the meiotic mechanism.



\section{Biological settings and extensions}
\label{sec:regimes}

\paragraph{Climate and repeated stress.}
A sustained or recurring environment may induce a similar meiotic map in many
generations. The immediate prediction is a reproducible change in $R_E$, $p_E$, or
both. Population-genomic effects require enough generations for the map history to
alter genealogies. Comparisons along thermal gradients should follow an experimental test of the
treatment-to-map relation.

\paragraph{Colonisation and range expansion.}
A bottleneck reduces standing variation. A recombination-map shift changes the
associations among the remaining alleles and the interference among selected loci.
It may reduce Hill--Robertson interference in some regions or break favourable and
locally adapted haplotypes in others. Its effect on adaptation depends on the
positions of the recombination changes relative to selected loci and haplotypes.

\paragraph{Parallel ecological transitions.}
Replicate populations in similar environments may show similar meiotic-state and
map responses. These responses may contribute to concordant patterns of diversity
or relative differentiation. Shared ancestral maps, structural variants, genetic
recombination modifiers, and parallel selection can produce similar patterns.
Support for the environmental-chromatin mechanism requires direct and repeated
observation of the meiotic-state and map changes.

\paragraph{Divergence with gene flow.}
Recombination affects linked selection and the persistence of multilocus
associations under migration. An environmental map shift can change the genomic background used to infer
barriers to gene flow. Peaks of $\FST$ have
several possible causes. A map shift can also weaken or strengthen adaptive
linkage, depending on its position. Joint models of migration, selection,
structural variation, and map history are required.

\section{Relation to existing theory and contribution}
\label{sec:novelty}

Separate parts of the model are established. Meiotic chromatin contributes to
recombination placement. Temperature can alter crossover number and distribution.
Linked selection relates diversity to recombination. Low within-population
diversity can raise relative differentiation. Changes in recombination maps can
alter background-selection signatures
\citep{underwood2018,lloyd2018,modliszewski2018,cruickshank2014,booker2022}.

Our contribution combines these results and makes four distinctions. First,
Eq.~\eqref{eq:decomp} separates total map-length change from spatial change while
allowing total recombination to vary. Second, Eq.~\eqref{eq:dp_general} gives an
exact sign rule for normalised placement, including regions with zero baseline
weight. Third, Eqs.~\eqref{eq:dlogpi_nonlocal} and \eqref{eq:pi_ratio_h} separate
the full nonlocal response from its local approximation. Fourth,
Eqs.~\eqref{eq:direct_memory_a}--\eqref{eq:direct_memory_m} distinguish direct
exposure, repeated induction, and inheritance. These distinctions define an
experimental sequence in which each link can be tested separately.

Booker et al.\ \citep{booker2022} model background selection under changing
recombination rates. The present model adds an environmentally responsive meiotic
input, the total--shape decomposition, and linked mechanistic and comparative
tests. It predicts changes in population-genomic variation when an environmental
cue alters recombination and the altered map recurs or persists over enough
generations.

\section{Assumptions and limits}
\label{sec:limits}

\paragraph{Chromatin accessibility.}
Accessibility may facilitate double-strand-break initiation. Crossover outcome
also depends on axis organisation, synapsis, sequence divergence, structural
heterozygosity, interference, and repair-pathway choice. The variable $q_E$ includes these contributions. Experiments can estimate it from
crossover data and the available molecular predictors.

\paragraph{Double-strand breaks and crossovers.}
A perturbation can increase break formation while leaving crossover number
unchanged. It can also redirect repair into noncrossovers or change crossover
maturation. Crossover maps provide the primary response variable. Break maps help
identify the stage at which the perturbation acts.

\paragraph{Crossover homeostasis.}
The obligate crossover and interference constrain crossover number. Temperature
experiments also show changes in $R_E$. Normalisation makes the shape component in
Eq.~\eqref{eq:decomp} integrate to zero. The total observed map can still change.

\paragraph{Local and nonlocal diversity responses.}
Equation~\eqref{eq:pi_ratio_h} gives a pointwise sign under the local
background-selection approximation. Equation~\eqref{eq:dlogpi_nonlocal} gives the full response. The distribution of
selected sites, chromosome boundaries, structural variants, and changes in
flanking regions can alter diversity at the site of interest.

\paragraph{Absolute divergence.}
The approximation $\Delta\dxy\approx0$ defines one regime. The analysis should
test it. Ancestral selection, mutation-rate variation, introgression, and
long-running map differences can change $\dxy$.

\paragraph{Recombination and adaptation.}
Higher recombination can separate favourable alleles from deleterious backgrounds
and reduce interference. It can also break coadapted or locally adapted haplotypes.
The model predicts a change in linkage structure. Genetic architecture and the
positions of selected loci determine the effect on adaptation.

\paragraph{Evidence for inheritance.}
Classify offspring effects according to which generations and germline precursors
received direct exposure. A test of transgenerational inheritance requires the
appropriate unexposed generations. The map-plasticity model still applies when the
chromatin state resets in every generation.

\section{Conclusion}

Studies of environmental effects on recombination should estimate total map length
and the spatial distribution of crossovers separately. The factorisation
$r_E=R_Ep_E$ and the decomposition
$\Delta r=\Delta R\pbar+\Rbar\Delta p$ provide this separation. A measured change
in crossover-placement weight gives an exact rule for the change in $p_E$.
Linked-selection theory then gives conditional predictions for diversity and
relative differentiation.

Tests should follow the same sequence as the model: measure meiotic state, estimate
crossover number and positions, establish the frequency or persistence of the map
change, and compare the resulting population-genomic patterns with simulations.
The formulation states the conditions under which an environmental effect on
meiotic chromatin can leave a population-genomic signal.

\bibliographystyle{plainnat}
\begin{thebibliography}{99}
\small

\bibitem[Baker et al.(2017)]{baker2017}
Baker Z, Schumer M, Haba Y, Bashkirova L, Holland C, Rosenthal GG, Przeworski M.
2017. Repeated losses of \PRDM-directed recombination despite the conservation of
\PRDM{} across vertebrates. \emph{eLife} 6:e24133.

\bibitem[Begun and Aquadro(1992)]{begun1992}
Begun DJ, Aquadro CF. 1992. Levels of naturally occurring DNA polymorphism
correlate with recombination rates in \emph{Drosophila melanogaster}.
\emph{Nature} 356:519--520.

\bibitem[Booker et al.(2022)]{booker2022}
Booker TR, Payseur BA, Tigano A. 2022. Background selection under evolving
recombination rates. \emph{Proceedings of the Royal Society B} 289:20220782.

\bibitem[Charlesworth et al.(1993)]{charlesworth1993}
Charlesworth B, Morgan MT, Charlesworth D. 1993. The effect of deleterious
mutations on neutral molecular variation. \emph{Genetics} 134:1289--1303.

\bibitem[Charlesworth(1998)]{charlesworth1998}
Charlesworth B. 1998. Measures of divergence between populations and the effect of
forces that reduce variability. \emph{Molecular Biology and Evolution} 15:538--543.

\bibitem[Choi et al.(2013)]{choi2013}
Choi K, Zhao X, Kelly KA, et al. 2013. Arabidopsis meiotic crossover hot spots
overlap with H2A.Z nucleosomes at gene promoters. \emph{Nature Genetics}
45:1327--1336.

\bibitem[Choi et al.(2018)]{choi2018}
Choi K, Zhao X, Tock AJ, et al. 2018. Nucleosomes and DNA methylation shape
meiotic DSB frequency in Arabidopsis transposons and gene regulatory regions.
\emph{Genome Research} 28:532--546.

\bibitem[Coop and Ralph(2012)]{coop2012}
Coop G, Ralph P. 2012. Patterns of neutral diversity under general models of
selective sweeps. \emph{Genetics} 192:205--224.

\bibitem[Coulton et al.(2020)]{coulton2020}
Coulton A, Burridge AJ, Edwards KJ. 2020. Examining the effects of temperature on
recombination in wheat. \emph{Frontiers in Plant Science} 11:230.

\bibitem[Cruickshank and Hahn(2014)]{cruickshank2014}
Cruickshank TE, Hahn MW. 2014. Reanalysis suggests that genomic islands of
speciation are due to reduced diversity, not reduced gene flow.
\emph{Molecular Ecology} 23:3133--3157.

\bibitem[De Storme and Geelen(2020)]{destorme2020}
De Storme N, Geelen D. 2020. High temperatures alter cross-over distribution and
induce male meiotic restitution in \emph{Arabidopsis thaliana}.
\emph{Communications Biology} 3:187.

\bibitem[Fernandes et al.(2024)]{fernandes2024}
Fernandes JB, Naish M, Lian Q, et al. 2024. Structural variation and DNA
methylation shape the centromere-proximal meiotic crossover landscape in
Arabidopsis. \emph{Genome Biology} 25:30.

\bibitem[Hill and Robertson(1966)]{hillrobertson1966}
Hill WG, Robertson A. 1966. The effect of linkage on limits to artificial
selection. \emph{Genetical Research} 8:269--294.

\bibitem[Hudson et al.(1992)]{hudson1992}
Hudson RR, Slatkin M, Maddison WP. 1992. Estimation of levels of gene flow from DNA
sequence data. \emph{Genetics} 132:583--589.

\bibitem[Hudson and Kaplan(1995)]{hudson1995}
Hudson RR, Kaplan NL. 1995. Deleterious background selection with recombination.
\emph{Genetics} 141:1605--1617.

\bibitem[Lloyd et al.(2018)]{lloyd2018}
Lloyd A, Morgan C, Franklin FCH, Bomblies K. 2018. Plasticity of meiotic
recombination rates in response to temperature in Arabidopsis.
\emph{Genetics} 208:1409--1420.

\bibitem[Martini et al.(2006)]{martini2006}
Martini E, Diaz RL, Hunter N, Keeney S. 2006. Crossover homeostasis in yeast
meiosis. \emph{Cell} 126:285--295.

\bibitem[Maynard Smith and Haigh(1974)]{maynard1974}
Maynard Smith J, Haigh J. 1974. The hitch-hiking effect of a favourable gene.
\emph{Genetical Research} 23:23--35.

\bibitem[Modliszewski et al.(2018)]{modliszewski2018}
Modliszewski JL, Wang H, Albright AR, et al. 2018. Elevated temperature increases
meiotic crossover frequency via the interfering (Type I) pathway in
\emph{Arabidopsis thaliana}. \emph{PLoS Genetics} 14:e1007384.

\bibitem[Nordborg et al.(1996)]{nordborg1996}
Nordborg M, Charlesworth B, Charlesworth D. 1996. The effect of recombination on
background selection. \emph{Genetical Research} 67:159--174.

\bibitem[Pan et al.(2011)]{pan2011}
Pan J, Sasaki M, Kniewel R, et al. 2011. A hierarchical combination of factors
shapes the genome-wide topography of yeast meiotic recombination initiation.
\emph{Cell} 144:719--731.

\bibitem[Son et al.(2025)]{son2025}
Son N, Kim H, Kim J, et al. 2025. The histone variant H2A.W restricts
heterochromatic crossovers in Arabidopsis. \emph{Proceedings of the National
Academy of Sciences USA} 122:e2413698122.

\bibitem[Szymanska-Lejman et al.(2026)]{szymanska2026}
Szymanska-Lejman M, Dziegielewski W, Wilhelm A, et al. 2026. Modifying meiotic
recombination by targeting chromatin regulators to crossover hotspots in
Arabidopsis. \emph{Science Advances} 12:eaeb2890.

\bibitem[Talbi et al.(2025)]{talbi2025}
Talbi M, Turner GF, Malinsky M. 2025. Rapid evolution of recombination landscapes
during the divergence of cichlid ecotypes in Lake Masoko. \emph{Evolution}
79:364--379.

\bibitem[Underwood et al.(2018)]{underwood2018}
Underwood CJ, Choi K, Lambing C, et al. 2018. Epigenetic activation of meiotic
recombination near \emph{Arabidopsis thaliana} centromeres via loss of H3K9me2 and
non-CG DNA methylation. \emph{Genome Research} 28:519--531.

\end{thebibliography}

\end{document}
